% Avsnitt 10.2-10.6, ur lectures/L10/appendix/a_can_frame_format.md.

\section{Introduktion till CAN}
\label{c10:sec:can}

CAN konstruerades för fordons- och industrisystem: många oberoende styrenheter (motor, bromsar,
instrumentpanel, sensorer) som utbyter korta meddelanden tillförlitligt, utan någon enskild
felpunkt.

Punkt-till-punkt-kablage slutar skala nästan omedelbart. Det ramade protokollet i
\secref{c10:sec:framing} förutsatte en enda delad länk; en styrenhet som pratar med en
instrumentpanel, en ABS-enhet och tre sensorer skulle behöva fem separata länkar, var och en med
eget kablage, egna kontakter och egen transceiver. En enda delad buss ersätter allihop: varje nod
kopplas in på samma ledningspar, och varje nod kan både lyssna och sända på den.

\subsection*{Den fysiska bussen}
En riktig CAN-buss är ett tvinnat differentialpar, \code{CAN_H} och \code{CAN_L}. En nod signalerar
en bit genom att driva isär de två ledningarna (\textbf{dominant}) eller låta dem flyta ihop genom
ett motstånd (\textbf{recessiv}). Ett externt transceiverchip sköter det analoga lagret.

Den här boken arbetar helt ovanför det: ett logiskt värde per bit, \code{0} för dominant och
\code{1} för recessiv. Det är exakt vad en transceiver presenterar för, och tar emot från, den
digitala logiken på vardera sidan.

\subsection*{Dominant och recessiv}
Hela arbitreringsschemat vilar på en enda egenskap:
\begin{itemize}
  \item \textbf{Recessiv (\code{1})} är bussens viloläge. Varje nod släpper bussen och
    transceivrarnas förspänningsnät drar båda ledningarna till samma nivå.
  \item \textbf{Dominant (\code{0})} hävdas genom att aktivt driva isär ledningarna. Om \emph{någon}
    nod driver dominant medan resten är recessiva läser bussen dominant, oavsett hur många noder som
    är inblandade eller vilken som hävdade det först.
\end{itemize}

Det är precis en \textbf{wired-AND}: betrakta varje nods utgång som en ingång till en busstäckande
AND-grind, och \code{0} dominerar. Riktiga transceivrar implementerar det i hårdvara med öppen
dränering; systemtestbänken i \chapref{c17:ch} modellerar identiskt beteende digitalt genom att
kombinera varje nods par \code{(tx_bus, bus_en)} på samma sätt.

\subsection*{Multimaster, och varför det kräver arbitrering}
De flesta inbyggda bussar besvarar frågan ”vem får sända?” med dekret: SPI har en enda master, och
I2C, som visserligen definierar en wired-AND-arbitrering ungefär som den nedan, körs nästan alltid
med en enda master den också. CAN gör multimaster till normalfallet: \textbf{vilken nod som helst
får sända så snart bussen är ledig}, och om två startar i samma ögonblick avgör bussen själv vem som
fortsätter. Det finns ingen omsändningsfördröjning och en kollision är inget fel alls. Det här är
CAN:s signaturidé, \textbf{arbitrering}.

Mekanismen i korthet; det genomräknade exemplet kommer i \secref{c10:sec:arbitration}:
\begin{enumerate}
  \item Varje meddelande börjar med sin identifierare, sänd bit för bit, MSB först.
  \item Varje sändande nod övervakar också bussen medan den sänder varje identifierarbit.
  \item Så länge biten den sände stämmer med det bussen läser fortsätter den sända.
  \item I det ögonblick en nod sänder recessivt (\code{1}) men observerar dominant (\code{0}) sänder
    någon annan nod en \code{0} där den själv sände en \code{1}. Den har förlorat arbitreringen: den
    slutar sända omedelbart och blir mottagare. Noden som fortsatte driva dominant vinner, utan att
    ha märkt någonting.
\end{enumerate}

Två följder av det:
\begin{itemize}
  \item \textbf{Lägre identifierare vinner alltid}, eftersom de bär fler inledande dominanta bitar.
    Därav ”prioriteter” i CAN-folkloren.
  \item \textbf{Varje identifierare måste vara unik.} CAN lägger det kravet på
    nätverkskonstruktören; kontrollern upprätthåller det inte i hårdvara.
\end{itemize}

\begin{aside}
\textbf{En förenkling att flagga tidigt.} Steg 4 beskriver en \emph{produktionsfärdig} kontroller,
och att sömlöst bli mottagare mitt i en ram är en del av varför riktig CAN behandlar en kollision
som en icke-händelse. Den här bokens \code{can_controller} (\chapref{c16:ch}--\ref{c17:ch}) slutar
driva i exakt samma ögonblick, vilket är den del som spelar roll för bussen, men överger sedan ramen
och sätter \code{error} i stället för att avkoda vinnarens meddelande. \Secref{c19:sec:gaps}
återkommer till det i sin lista över skillnader mot en produktionsfärdig kontroller.
\end{aside}

\Canref{1} tar bussen på egen hand: topologin, den dominanta och recessiva asymmetrin,
wired-AND:et, de två termineringsmotstånden, och sambandet mellan bithastighet och busslängd
\textemdash{} 1~Mbit/s räcker ungefär 40 m, och skälet är arbitreringen, som kräver att en bit hinner
ut till bussens bortersta nod och tillbaka innan biten är slut.

% ------------------------------------------------------------------------------------------
\section{Broadcast, inte adressering}
\label{c10:sec:broadcast}

Det ramade protokollet i \secref{c10:sec:framing} behövde DST och SRC eftersom fler än en nod kunde
lyssna. CAN står inför samma problem på samma sorts buss och besvarar det helt annorlunda: \textbf{en
CAN-ram bär ingen destinationsadress alls.}

Varje nod ser varje ram. Identifieraren som vinner arbitreringen namnger ingen mottagare; den
namnger \emph{sortens meddelande}, utsänt till alla och relevant för den som bryr sig. Varje nod
avgör själv om en identifierare betyder något, typiskt genom att jämföra den mot en liten uppsättning
den är intresserad av, i mjukvara eller i kontrollerns filtreringshårdvara. En nod som inte är
intresserad kastar helt enkelt ramen; ingenting på protokollnivå styrde undan den.

\paragraph{Exempel}
Anta att en konstruktion reserverar identifierare enligt konvention,
$\mathit{StatusReq} = \code{0x300} + \mathit{nodeId}$ och
$\mathit{StatusResp} = \code{0x400} + \mathit{nodeId}$. Sensornod 3 tillfrågas om status på
\code{0x303} och svarar på \code{0x403}. Varje nod tar emot ramen \code{0x303}; bara den som känner
igen den som sin egen förfrågningsidentifierare svarar. Inget fält sa ”det här är till nod 3”.

Det är ett genuint annorlunda designval. Protokollet i \secref{c10:sec:framing} skalar adressering
genom att bredda DST och SRC när antalet noder växer; CAN skalar genom konvention för hur
identifierare tilldelas och filtreras, och det finns inget adressfält att bredda.

% ------------------------------------------------------------------------------------------
\section{CAN-ramformatet}
\label{c10:sec:frameformat}

\lecfig[0.86]{lectures/L10/appendix/images/standard_can_frame_format.png}{CAN:s standarddataram,
fält för fält, i sändningsordning.}{c10:fig:canframe}

Den här boken implementerar CAN:s \textbf{standarddataram (basformatet)}, bit för bit som standarden
definierar den: en 11-bitars identifierare och upp till 8 databytes. Fjärramar, utökade 29-bitars
identifierare och allt standarden definierar \emph{runt} dataramar (fel- och överlastramar, och den
3-bitars intermission som skiljer en ram från nästa) ligger utanför bokens ram, men varje ram den
här kontrollern lägger ut på ledningen är en fullt standardenlig dataram. Det är samma förenkling
som riktiga konstruktioner ofta utgår från, och den räcker gott och väl för att bygga en genuin,
fungerande kontroller.

Diagrammet ger sändningsordningen och varje fälts bredd. Bitstoppning gäller varje fält från
\textbf{SOF till och med CRC-fältet}, och inte CRC-avgränsaren, ACK-fältet eller EOF.

\subsection*{Fält för fält}
\paragraph{SOF, 1 bit, alltid dominant}
Sänds ut på en i övrigt ledig, recessiv buss, så den producerar den enda
recessiv-till-dominant-flank varje nod väntar på. Den flanken markerar både att en ram börjar och
resynkroniserar varje nods bittajming mot den; \code{bit_timer.resync} i \chapref{c12:ch} är precis
det här.

\paragraph{Identifierare, 11 bitar, MSB först}
Fungerar samtidigt som meddelandets arbitreringsprioritet. Varje sändande nod jämför varje bit den
sänder mot vad bussen faktiskt läser, och faller ifrån i det ögonblick den sänder recessivt men
observerar dominant.

\paragraph{RTR, 1 bit, alltid dominant här}
Remote Transmission Request-biten avslutar det 12 bitar långa arbitreringsfältet: dominant markerar
en dataram, recessiv en \emph{fjärram}, en begäran om att den som äger identifieraren sänder sin
data. Den här boken stöder bara dataramar, så dess kontroller sänder alltid RTR dominant, och en
mottagare som samplar den recessiv överger ramen och sätter \code{error} (\chapref{c17:ch}
implementerar den kontrollen). Formellt spänner arbitreringen över identifieraren \emph{och} RTR; i
riktig CAN är det det som låter en dataram slå en fjärram med samma identifierare. När varje nod här
sänder RTR dominant avgör identifieraren alltid ensam.

\paragraph{Kontrollfält, 6 bitar: IDE, r0, DLC}
\code{IDE} är alltid \code{0} här; \code{1} skulle markera en ram med utökad identifierare, utanför
bokens ram. \code{r0} är reserverad, alltid \code{0}. \code{DLC} är det 4 bitar breda längdfältet,
\code{0}--\code{8} bytes kodade som \code{0000}--\code{1000}. Det är längdfältet från
\secref{c10:sec:framing} i ny skepnad: utan det skulle en mottagare inte kunna avgöra om datafältet
slutar efter den här byten eller fortsätter.

\paragraph{Datafält, 0 till 8 bytes}
Sänds MSB först, byte 0 först, med den längd DLC annonserade. När DLC är \code{0} saknas fältet helt
och CRC:n följer direkt på kontrollfältet.

\paragraph{CRC-fält, 15 bitar plus en avgränsare}
CRC:n täcker varje bit från SOF till och med datafältets slut, och är långt mer robust än den
summerade checksumman i \secref{c10:sec:framing}. En detalj spelar roll längre fram: den täcker bara
de \emph{riktiga} bitarna. Stoppbitar sätts in efter att CRC:n beräknats på vägen ut och kastas innan
den kontrolleras på vägen in, så de kommer aldrig in i beräkningen (\chapref{c14:ch} och
\ref{c15:ch} är där den grindningen byggs). \Chapref{c13:ch}:s motor beräknar den bitseriellt, en bit
i taget medan ramen sänds eller tas emot, utan något separat ”beräkna checksumman”-steg på slutet.

\textbf{CRC-avgränsaren} är den 16:e biten, alltid recessiv och aldrig stoppad. Den finns för att en
bit som är oberoende av stoppbitsregeln alltid ska följa på CRC-värdet, vilket är det som låter en
mottagare hitta gränsen mellan CRC- och ACK-fälten.

\paragraph{ACK-fält, 2 bitar}
Sändaren sänder båda recessiva. Varje \emph{annan} nod som tog emot ramen korrekt (ramningen intakt,
stoppningen respekterad, CRC:n giltig) drar ACK-luckan dominant. Sändaren kontrollerar aldrig sin
egen CRC mot sig själv, så den dominanta biten från någon annan är dess enda bekräftelse på att ramen
kom fram intakt någonstans. \textbf{ACK-avgränsaren} är recessiv och ostoppad, och markerar gränsen
av samma skäl som CRC-avgränsaren.

\paragraph{EOF, 7 recessiva bitar, ostoppade}
Sju identiska bitar i rad är långt fler än stoppbitsregeln någonsin skulle tillåta inuti ett stoppat
fält, vilket är precis det som gör EOF omisskännligt. Det håller bara om varje nod är överens om att
inte stoppa det, vilket är varför EOF och de två avgränsarna definieras som undantagna från
stoppning.

\Canref{2} går igenom samma ram fält för fält, och tar med det som ligger utanför den här bokens
kontroller: fjärramarnas RTR-semantik, vad ACK-luckan betyder för en sändare som läser tillbaka den,
och de tre bitars intermission som skiljer en ram från nästa.

% ------------------------------------------------------------------------------------------
\section{Bitstoppning}
\label{c10:sec:stuffing}

\textbf{Regeln.} Genom varje stoppat fält, SOF till och med CRC-fältet: efter \textbf{fem bitar i
följd med samma värde} sätter sändaren in en extra bit med \emph{motsatt} värde. Den biten bär ingen
data, och varje mottagare kastar den genom att räkna på samma sätt.

\textbf{Varför den finns}, av två besläktade skäl:
\begin{enumerate}
  \item \textbf{Klockåtervinning.} Varje nods bittimer går på sin egen oscillator. En lång följd av
    identiska bitar innehåller inga övergångar, så en mottagare har ingenting som bekräftar att den
    fortfarande ligger i linje med sändarens bitgränser. Att garantera en övergång minst var femte
    riktig bit begränsar hur långt dess tajming kan driva. Riktig CAN resynkroniserar på de flanker
    bitstoppningen garanterar; den här bokens kontroller är enklare och resynkroniserar bara vid SOF
    (\code{bit_timer.resync}, \chapref{c12:ch}), vilket gör den begränsade driften ännu viktigare.
  \item \textbf{Att skilja riktig data från EOF.} EOF är sju identiska recessiva bitar, ostoppade.
    Ett stoppat fält kan aldrig av misstag producera den följden, vilket är det som gör sju
    ostoppade recessiva bitar till ett otvetydigt ramslut snarare än något identifieraren eller
    datafältet kunde ha producerat av en slump.
\end{enumerate}

\paragraph{Genomräknat exempel}
För att sända den 8 bitar långa sekvensen \code{11111000} slår regeln till efter den femte
\code{1}:an i följd: sätt in en \code{0}, fortsätt sedan med den riktiga datan. Varje kolumn nedan
är en bit som faktiskt ligger på ledningen, så raden med riktiga bitar lämnar en lucka där
stoppbiten sitter:

\begin{textdrawing}
Riktiga bitar:            1  1  1  1  1  -  0  0  0
Sänd bitström:            1  1  1  1  1  0* 0  0  0
Antal i följd:            1  2  3  4  5  1  2  3  4
                                         ^ stoppbit - inte riktig data
\end{textdrawing}

Läs antalsraden mot den sända raden, inte mot raden med riktiga bitar: regeln räknar bitar på
ledningen. Stoppbiten är själv den första biten i den nya följden, vilket är varför antalet står som
\code{1} under den och \code{2} under den första riktiga \code{0} som följer. \Chapref{c14:ch} bygger
precis den här räknaren i hårdvara.

En mottagare räknar samma fem \code{1}:or, förväntar sig att nästa bit är en stoppbit, och kastar
den om den inverterar följden. Dyker det i stället upp en \emph{sjätte} \code{1} i följd där
stoppbiten skulle ha suttit är det en stoppningsöverträdelse: ett detekterat ramfel, behandlat som
fatalt, rapporterat av \code{rx_shift_reg.stuff_error} i \chapref{c15:ch}.

\begin{warning}
\textbf{En finess värd att nämna redan nu, eftersom den överraskar folk senare.} Räkningen av bitar
i följd nollställs \emph{inte} vid fältgränser. En följd som börjar i ett fält och fortsätter in i
nästa räknas rakt igenom, och kan utlösa en stoppbit precis vid gränsen. Stoppningen bryr sig bara
om bitarna som faktiskt ligger på ledningen, aldrig om vilket namngivet fält de tillhör.
\end{warning}

\Canref{3} härleder femregeln ur klockåtervinningsproblemet och tar CRC-15:an i samma kapitel, av
det skäl \chapref{c13:ch} och \ref{c14:ch} sedan lever med: stoppningen och checksumman delar
bitström men inte bitar, och den som förstår den ena utan den andra får fel CRC.

% ------------------------------------------------------------------------------------------
\section{Arbitrering, genomräknad i sin helhet}
\label{c10:sec:arbitration}

Noderna P (identifierare \code{0x0F0}) och Q (identifierare \code{0x100}) börjar sända vid samma
SOF. Båda sänder SOF som \code{0}, så ingen strid uppstår än. Sedan identifierarna, MSB först:

\begin{textdrawing}
             bit:  1  2  3  4  5  6  7  8  9 10 11
P (0x0F0) sänder:  0  0  0  1  1  1  1  0  0  0  0
Q (0x100) sänder:  0  0  1  -  -  -  -  -  -  -  -   (förlorar vid bit 3, slutar driva)
bussen läser:      0  0  0  1  1  1  1  0  0  0  0
\end{textdrawing}

Bitarna 1--2 stämmer för båda noderna, så ingenting händer. Vid bit 3 sänder P dominant och Q
recessivt, så bussen läser dominant. Q, som övervakar medan den sänder, ser sin egen \code{1} mot en
buss som läser \code{0}: \textbf{Q har förlorat arbitreringen} och slutar driva omedelbart. P ser
\code{0} på en buss som läser \code{0}, exakt vad den förväntade sig, och fortsätter utan att
någonsin få veta att en strid ägde rum.

I en produktionsfärdig kontroller skulle Q nu ta emot resten av P:s ram. Den här bokens
\code{can_controller} slutar driva i samma ögonblick men överger ramen, som flaggades ovan och som
\secref{c19:sec:gaps} återkommer till.

Mekanismen generaliserar till godtyckligt antal samtidiga sändare: den numeriskt lägsta
identifieraren vinner (likvärdigt: den längsta följden av inledande dominanta bitar), och varje
förlorare faller ifrån i det ögonblick dess recessiva bit övertrumfas, aldrig senare än exakt den
biten.

\Canref{4} räknar igenom både tvånods- och trenodsfallet, och tar sedan de två frågor det här
kapitlet inte ställer: vad svält och prioritetsinversion betyder på en buss där låga identifierare
alltid vinner, och vad en sändare aldrig får veta om varför den förlorade.

% ------------------------------------------------------------------------------------------
\section{Bittajming och sampelpunkten}
\label{c10:sec:bittiming}

Arbitrering fungerar bara om varje nod är överens om \emph{när} en bit ligger på bussen, vilket
väcker en fråga: vid vilken punkt inom en bitperiod läser en nod den?

Varje nod delar upp sin egen klocka i bitperioder med fast längd med hjälp av sin egen lokala
bittimer. Eftersom de oscillatorerna är oberoende, och eftersom signaler tar tid på sig att utbreda
sig längs bussen, samplar en nod \textbf{inte} vid periodens början. Den väntar till en
\textbf{sampelpunkt} längre in i perioden, då den sända nivån hunnit stabilisera sig.

Den här boken samplar 70~\% in i varje bitperiod, ett rimligt och enkelt val för en kort buss på ett
enda kort vid 1~Mbit/s.

Riktiga CAN-kontrollrar gör sampelpunkten konfigurerbar och delar upp bitperioden i namngivna
segment (utbredningssegment, fassegment) avstämda mot busslängd, utbredningsfördröjning och
oscillatortolerans. Det ligger utanför den här bokens fasta sampelpunkt, men inte utanför
\canref{5}, som delar upp bitperioden i sina segment, räknar ut tajmingen ur en klockfrekvens och en
bithastighet, och beskriver den resynkronisering vid varje flank som den här kontrollern ersätter
med en enda puls vid SOF.

\Secref{c10:sec:candef} gör om de 70~\% till den konkreta VHDL-konstanten \code{SAMPLE_TICK}, härledd
ur klockfrekvensen och bithastigheten i stället för hårdkodad.
